% Apoyo visual del resultado ya deducido para la componente (b,c,d)=(t,t,r).
% La normalización elimina los parámetros sin asignarles valores numéricos.
\begin{frame}[label=nav-frame-121]{Lectura radial de la componente calculada}
\small
\[
C(r)\equiv\left|\nabla_aP_\beta^{attr}\right|
=\frac{\ell^2|\beta_0|p^2}{4r^3},\qquad r>0.
\]
\begin{columns}[c,onlytextwidth]
\begin{column}{.58\textwidth}
\centering
\begin{tikzpicture}[x=1.5cm,y=3.3cm,font=\footnotesize]
  \foreach \yt/\lab in {0/0,0.25/{0{,}25},0.5/{0{,}5},0.75/{0{,}75},1/1}{
    \draw[black!12] (1,\yt) -- (4.1,\yt);
    \node[anchor=east] at (.94,\yt) {$\lab$};
  }
  \draw[black!65] (1,0) -- (4.15,0);
  \draw[black!65] (1,0) -- (1,1.05);
  \foreach \xt in {1,2,3,4}{
    \draw[black!65] (\xt,0) -- (\xt,-.02);
    \node[anchor=north] at (\xt,-.03) {$\xt$};
  }
  \draw[structure.fg,line width=1.2pt,domain=1:4,samples=121,smooth,variable=\x]
    plot (\x,{1/(\x*\x*\x)});
  \fill[structure.fg] (2,.125) circle[radius=1.6pt];
  \node[anchor=south west,text=structure.fg] at (2.03,.14) {$1/8$};
  \fill[structure.fg] (4,.015625) circle[radius=1.6pt];
  \node[anchor=south east,text=structure.fg] at (4,.055) {$1/64$};
  \node at (2.55,-.19) {$r/r_0$};
  \node[rotate=90] at (.35,.52) {$C(r)/C(r_0)$};
\end{tikzpicture}
\end{column}
\begin{column}{.39\textwidth}
Para $\beta_0p\neq0$, elegimos un radio de referencia $r_0>0$:
\[
\frac{C(r)}{C(r_0)}=\left(\frac{r_0}{r}\right)^3.
\]
Al duplicar el radio, la componente se reduce a $1/8$ de su valor.

\medskip
A radio finito sigue siendo no nula. Su decaimiento asintótico
no establece la condición LL.
\end{column}
\end{columns}
\vspace{2mm}
\footnotesize
Valor absoluto de la componente calculada en el ansatz, no una norma
tensorial. Los casos $\beta_0=0$ o $p=0$ anulan la componente y quedan
fuera de esta normalización.
\end{frame}
